从单变量走到多变量，表面上只是把自变量从一个数 \(x\) 换成一个向量 \(\mathbf x\)，但脚下的几何图景已经完全不同了。

在单变量微积分里，你只能在一根数轴上向左或向右挪动，导数就是一个斜率。到了多元的情形，你站在一片高低起伏的地形上：
\begin{itemize}
    \item 脚下有无穷多个方向可以迈步，每个方向的陡峭程度（方向导数）都不一样；
    \item 局部的线性近似不再是一个数，而是一张切平面（一个线性映射）；
    \item 积分也不再只是曲线下的面积，而是发生在区域、曲线和曲面上的累积。
\end{itemize}

这些升级不是学究式的推演，它们直接构成了现代优化与机器学习的几何基石。读本单元的九篇时，可以对照下面几条主线：

\begin{enumerate}
    \item \textbf{为什么偏导数不够用？}\par
    只沿着坐标轴切片观察（偏导数），可能会完全错过地形在对角线方向的剧烈变化。只有当``切平面''存在时（可微性），局部线性近似才对所有方向同时有效——正因如此，梯度才配被称为最速上升方向。

    \item \textbf{梯度是如何被``尺子''决定的？}\par
    ``最陡的下山方向''并不是偏导数的简单堆砌，它取决于你如何定义``一步''的距离（内积）。换用不同的度量，最速方向就会改变——逐坐标自适应（如 Adam）、预条件与自然梯度的核心区别，全在于选择了怎样的空间尺子。

    \item \textbf{链式法则的几何本质}\par
    高维函数的复合，本质上是线性映射的复合。矩阵乘法只是它在坐标系下的落地形式；而在计算图里选择先乘左边还是先乘右边（结合律），就直接决定了前向模式与反向传播的计算效率差距。

    \item \textbf{高维地形里的鞍点}\par
    在高维空间里，严格的山顶和山谷（极值点）非常罕见，梯度归零的地方绝大多数是``鞍点''——沿一个方向看是山谷，沿另一个方向看却是山脊。要区分它们，得看 Hessian 在各个方向上的曲率（特征值）。

    \item \textbf{约束与惩罚的等价性}\par
    限制在曲面上行走（约束优化），相当于把``所有方向都平坦''削弱为``只在切向平坦''。最优点处目标梯度与约束梯度必定共线，而 Lagrange 乘子正是约束条件的``影子价格''。

    \item \textbf{空间扭曲与体积伸缩}\par
    坐标变换不仅改变位置，还会拉伸空间。Jacobian 行列式正是这个局部体积的伸缩因子——概率密度变换时需要除以它，生成流模型（Flow-based Models）计算对数似然时累加的就是它的对数。
\end{enumerate}

后三篇转向累积。无论是区域上的多重积分、曲线上的做功，还是曲面上的通量，最后都汇聚到同一个事实上：区域内部无数微小的局部变化累积起来，恰好等于边界上观察到的净效应。Green、Gauss 和 Stokes 公式，就是这句话在不同维度下的写法。

\subsection*{怎么读这一章}

\begin{itemize}
    \item \textbf{主线（前六篇）}：建议按顺序阅读，它们共同回答一个核心问题——\emph{多元函数在一点附近究竟怎样变化？} 读到能够直观说清偏导与可微的差别、梯度方向为什么依赖内积、Hessian 特征值代表什么，这一趟的收获就足够了。
    \item \textbf{支线（后三篇）}：若你的重点在概率分布与生成模型，第七篇（变量替换）是必读内容；若你需要处理物理场、流体或连续介质，第八、九篇则是核心主线。复杂曲面上的具体计算不必死磕，用到时再回来查阅。
\end{itemize}

\subsection*{本章篇目}

\begin{itemize}
    \item \hyperref[sec:02-Calculus/06-Multivariable-Calculus/01]{``多元函数、极限与连续''}：为什么沿着几条直线检查平滑，从来不构成全局连续的证明；
    \item \hyperref[sec:02-Calculus/06-Multivariable-Calculus/02]{``偏导数为何不等于可微''}：数值梯度检查全部通过，为什么不代表局部存在合格的线性模型；
    \item \hyperref[sec:02-Calculus/06-Multivariable-Calculus/03]{``梯度与方向导数''}：负梯度到底凭什么被称为``最速下降方向''？
    \item \hyperref[sec:02-Calculus/06-Multivariable-Calculus/04]{``Jacobian 与多元链式法则''}：反向传播的全部数学图景，以及它为何比前向计算更便宜；
    \item \hyperref[sec:02-Calculus/06-Multivariable-Calculus/05]{``Hessian 与无约束极值''}：当梯度归零时，靠特征值区分山谷、山顶与马鞍；
    \item \hyperref[sec:02-Calculus/06-Multivariable-Calculus/06]{``约束几何与 Lagrange 乘子''}：硬性约束与软性惩罚项，为何是同一件事的两种写法；
    \item \hyperref[sec:02-Calculus/06-Multivariable-Calculus/07]{``多重积分与变量替换''}：扭曲坐标系时空间如何缩放，流模型的对数行列式由此而来；
    \item \hyperref[sec:02-Calculus/06-Multivariable-Calculus/08]{``向量场、曲线积分与势函数''}：沿着轨迹走，什么时候能积累能量，什么时候只会原地绕圈；
    \item \hyperref[sec:02-Calculus/06-Multivariable-Calculus/09]{``通量、散度、旋度与边界定理''}：内部变化的累积，如何精确映射为边界上的总效应。
\end{itemize}

这九篇中建立的几何直觉——梯度、Jacobian、Hessian、乘子与行列式——将在后续的凸优化、矩阵微积分和信息几何中反复出现。比计算技巧更重要的，是带着这些对象在思维中建立起清晰的图景。